\documentclass[11pt,a4paper]{article}
\usepackage[margin=1in]{geometry}
\usepackage[T1]{fontenc}
\usepackage{graphicx}
\usepackage{booktabs}
\usepackage{float}
\usepackage{hyperref}
\usepackage{caption}
\usepackage{enumitem}
\usepackage{xurl}
\usepackage{array}
\emergencystretch=3em

% Figure if the file exists, otherwise a visible placeholder box.
\newcommand{\safefig}[3]{%
  \begin{figure}[H]\centering
  \IfFileExists{../reports/figures/#1}{\includegraphics[width=#2\linewidth]{../reports/figures/#1}}{%
    \fbox{\parbox{0.8\linewidth}{\centering\vspace{1.2cm}\detokenize{#1} not found: run the EDA scripts first.\vspace{1.2cm}}}}%
  \caption{#3}\end{figure}}

\title{Federated Learning for EEG-Based Working-Memory Decoding\\
\large Phase 1 and Phase 2 Report (Deliverable 1)}
\author{Eman Adil \and Dua Sarfraz \and Ayza Ahmad\\[3pt] {Artificial Intelligence}}
\date{\today}

\begin{document}
\maketitle

\begin{abstract}
This report presents the problem definition and data foundations of a pilot study on federated learning (FL)
for decoding working-memory load from EEG. We use OpenNeuro dataset \texttt{ds003655} (verbal working-memory
task); two participants have been acquired so far. We state the research problem and justify the use of machine
learning (Phase 1), then document data acquisition, validation, a 240-trial table, signal-quality investigation,
first-pass spectral features (theta, alpha, beta band power for 19 EEG channels), the cleaning and scaling
strategy, and exploratory data analysis (Phase 2). Because only two participants are available, this is a pilot
and no claims are made about generalization. Several signal-quality and label-validity questions remain open and
are reported explicitly rather than hidden.
\end{abstract}

\tableofcontents
\newpage

%===========================================================
\section{Phase 1: Problem Definition}

\subsection{Problem Statement}
Working memory is the short-term maintenance and manipulation of information, and it is central to reasoning and
learning. EEG measures brain activity non-invasively and cheaply, so decoding working-memory demand from EEG is of
interest for cognitive-state monitoring and brain--computer interfaces. Two practical obstacles limit this:
\begin{enumerate}[leftmargin=*]
  \item \textbf{Inter-subject variability.} EEG differs strongly between people (anatomy, electrode contact, noise,
  individual rhythm frequencies), so a model trained on some people often degrades on others, and data from many
  participants is needed.
  \item \textbf{Privacy and data silos.} EEG is sensitive personal data. Pooling recordings from different labs or
  clinics is often restricted by consent and regulation.
\end{enumerate}
Federated learning trains a shared model by exchanging model updates instead of raw recordings, which addresses the
second obstacle. However, each participant's data is distributed differently (non-IID), which is exactly the first
obstacle and is known to make federated optimisation harder \cite{mcmahan2017}. This project studies that
interaction for EEG working-memory decoding.

\subsection{Research Questions}
\begin{enumerate}[leftmargin=*]
  \item How does a federated model compare with a centralized baseline for working-memory condition decoding?
  \item How sensitive are results to participant-level distribution differences?
  \item Does a clearly defined personalization strategy improve cross-subject generalization?
\end{enumerate}
These are planned experiments, not results at this stage.

\subsection{Target Audience and Research Gap}
This is an original research (R\&D) project. The intended audience is researchers in neuro-engineering,
brain--computer interfaces and privacy-preserving machine learning. The gap is a reproducible, small-scale comparison
of centralized, federated and personalized learning for EEG working-memory decoding under non-IID, cross-subject
conditions, in which data quality and label validity are examined rather than assumed.

\textbf{Scope.} Only two participants are available. Findings describe a pipeline and its behaviour on two people and
cannot support claims about the wider population.

\subsection{Justification for Machine Learning}
A rule-based program would need hand-written rules from EEG measurements to working-memory demand. This is unsuitable
because:
\begin{itemize}[leftmargin=*]
  \item \textbf{High-dimensional, weak, distributed signal.} Each trial is a 19-channel time series at 500\,Hz;
  memory-load effects appear as subtle spectral changes spread over channels, with no single threshold separating
  conditions.
  \item \textbf{Subject-specific patterns.} The same state looks different across people, so fixed rules do not transfer,
  whereas learned models can adapt to these differences.
  \item \textbf{Low signal-to-noise ratio.} Artifacts (eye movement, muscle activity, amplifier saturation; see
  Section~\ref{sec:quality}) require statistical models that weigh evidence over many features.
  \item \textbf{The research question itself is about learning.} Federated optimisation and personalization are
  properties of learning algorithms and have no rule-based counterpart.
\end{itemize}

%===========================================================
\section{Phase 2: Data Acquisition}

\subsection{Dataset and Collection}
The data are from OpenNeuro dataset \texttt{ds003655}, task \texttt{VerbalWorkingMemory}
(\url{https://openneuro.org/datasets/ds003655}), a modified Sternberg task with retention (R) and manipulation (M)
conditions at set sizes 5, 6 and 7, giving the conditions 5R, 6R, 7R, 5M, 6M and 7M. Two participants,
\texttt{sub-001} and \texttt{sub-002}, were downloaded from the public repository. The local inventory
(\texttt{reports/download\_inventory.csv}, produced by \texttt{01\_check\_downloads.py}) records eight files per
participant and about 269.05\,MB in total. Raw files are excluded from Git through \texttt{.gitignore}. No download
API log was captured.

\begin{table}[H]\centering
\caption{Recording metadata.}
\begin{tabular}{lcc}
\toprule
Property & sub-001 & sub-002 \\
\midrule
Sampling rate & 500 Hz & 500 Hz \\
Channels recorded & 21 & 21 \\
EEG channels & 19 & 19 \\
EOG channels & EOGv, EOGh & EOGv, EOGh \\
Duration & 3077 s & 3035 s \\
\bottomrule
\end{tabular}
\end{table}
The EEG channels are Fp1, Fp2, F7, F3, Fz, F4, F8, T3, C3, Cz, C4, T4, T5, P3, Pz, P4, T6, O1, O2.

\subsection{Trial Structure}
\texttt{04\_build\_trial\_table.py} builds \texttt{reports/trial\_table.csv} from the event markers: 240 trials
(120 per participant), 20 per condition per participant. Each row holds the onsets of the baseline, encoding, delay
and probe phases, the response information and quality flags. The delay-to-probe interval is about 6.7\,s, but all
scripts use the actual event onsets.

%===========================================================
\section{Phase 2: Validation and Data-Centric Cleaning}
\label{sec:quality}

\subsection{File and Event Validation}
Both recordings loaded correctly, with 120 trial starts and 20 baseline, encoding, delay and probe markers per
condition. The event tables had no exact duplicate rows and no onsets outside the recording. These checks confirm
structure, not scientific validity.

\subsection{Response-Label Uncertainty}
There are 244 response markers for 240 probes; some are late, duplicated or ambiguous. The trial table flags eight trials
for response ambiguity (6 in sub-001, 2 in sub-002), and the later probe-based audit marks 225 of 240 rows as PASS and 15
as REVIEW; the two counts come from different checks and are to be reconciled. Late responses are judged with a five-second heuristic that is not a
validated rule. Uncertain cases are kept, not silently resolved, and the event-to-label mapping is provisional; no
supervised model will be trained before the labels are validated.

\subsection{Channel-Type Correction}
MNE initially labelled \texttt{EOGv} and \texttt{EOGh} as EEG. They are now set to type \texttt{eog} before channel selection,
so features use only the 19 EEG channels (\texttt{12\_inspect\_channel\_types.py}, \texttt{18\_extract\_features.py}).

\subsection{Signal-Quality Findings}
\textbf{sub-001 boundary values.} In all 21 channels the minimum value reaches about $-10{,}239.69\,\mu$V (maximum
absolute value 10{,}239.69\,$\mu$V). Between 332.5 and 333.5\,s each channel has many samples near this lower boundary,
including consecutive runs of about 0.48--0.52\,s, and the same boundary was not observed in sub-002 in that interval.
This is suspicious for clipping, saturation or a data-representation issue, but \textbf{the cause is not confirmed}
(\texttt{reports/signal\_quality\_decisions.md}).

\textbf{Amplitude screening.} Each trial window is flagged if its maximum channel peak-to-peak or absolute amplitude
exceeds a heuristic threshold of 500\,$\mu$V.
\begin{table}[H]\centering
\caption{Heuristic amplitude flags per trial window (not validated rejection rules).}
\begin{tabular}{lccc}
\toprule
Participant & No flag & Peak-to-peak only & Both flags \\
\midrule
sub-001 & 112 & 6 & 2 \\
sub-002 & 0 & 0 & 120 \\
\midrule
Total & 112 & 6 & 122 \\
\bottomrule
\end{tabular}
\end{table}
\begin{table}[H]\centering
\caption{Window amplitude by flag group ($\mu$V).}
\begin{tabular}{lccc}
\toprule
Group & Median p--p & Mean p--p & Median abs.\ amp. \\
\midrule
No flag & 406.66 & 411.35 & 258.31 \\
Peak-to-peak only & 524.28 & 550.06 & 371.22 \\
Both flags & 2550.53 & 2588.39 & 1325.66 \\
\bottomrule
\end{tabular}
\end{table}
Every sub-002 trial is flagged, which does not by itself mean every trial should be discarded. Whole-record
channel ranges in sub-002 are large but vary by channel (e.g.\ Pz 5461.00, T3 5318.31, O1 5506.38\,$\mu$V), and the EOG
channels also have large extrema. To tell a single bad channel from a participant-wide difference, the per-channel,
per-trial audit (\texttt{24\_inspect\_trial\_channel\_amplitudes.py}, output
\texttt{reports/trial\_channel\_amplitude\_audit.csv}) was computed for every delay-to-probe window
(Table~\ref{tab:chanaudit}).
\begin{table}[H]\centering
\caption{Per-trial channel amplitude audit by participant (120 windows each; $\mu$V).}\label{tab:chanaudit}
\begin{tabular}{>{\raggedright\arraybackslash}p{5.4cm}>{\raggedright\arraybackslash}p{4.1cm}>{\raggedright\arraybackslash}p{4.1cm}}
\toprule
Quantity & sub-001 & sub-002 \\
\midrule
Max channel peak-to-peak, median (range) & 410.8 (317.4--1612.2) & 2551.4 (2444.9--2989.5) \\
Max absolute amplitude, median & 260.5 & 1325.7 \\
Median-over-channels peak-to-peak, mean & 256.5 & 613.0 \\
Channels above 500\,$\mu$V p--p per window, mean (range) & 0.24 (0--19) & 13.2 (12--19) \\
Channel with the largest p--p & Fp2 50, T3 45, O2 24, F8 1 trials & T3 in all 120 trials \\
\bottomrule
\end{tabular}
\end{table}
Two observations follow. First, the sub-002 difference is not confined to one channel: the typical channel is itself
about 2.4 times larger than in sub-001, and 12 or more of the 19 EEG channels exceed 500\,$\mu$V in every window. Second,
T3 has the largest range in every sub-002 window, so a channel-specific problem may be superimposed on the general
difference. In sub-001 the flags are rare and concentrated in a few windows (the maximum of 19 channels over threshold
belongs to a single window). The cause of the sub-002 pattern (recording or reference differences, scaling, drift, or
genuine activity) is \textbf{still not established}. Note also that even the unflagged sub-001 windows have a median
channel peak-to-peak range of roughly 250\,$\mu$V, which is large for EEG; we hypothesise that this reflects the
absence of high-pass filtering, but this has not been tested.

\textbf{Decision policy.} No trial or participant has been excluded, and raw files are never modified. Rejection or repair
rules will be chosen only after this investigation and recorded in \texttt{signal\_quality\_decisions.md}. The
\texttt{overlaps\_known\_suspicious\_interval} column currently reports no overlap with feature windows; this still has
to be re-checked against actual onsets.

\subsection{Tabular Cleaning of the Feature Table}
Continuous EEG is not cleaned as if extreme samples were ordinary tabular outliers. For the derived trial-level feature
table, \texttt{21\_clean\_tabular\_features.py} (i) reports and removes exact duplicate rows, (ii) imputes missing numeric
values with the median and missing categories with ``Unknown'', (iii) flags IQR outliers ($1.5\times$IQR) without removing
them, and (iv) applies standard scaling. Run on the 240-row feature table, it found 0 exact duplicate rows and 0
missing values in all 60 numeric columns (57 band-power features, window duration and the two amplitude columns). The IQR
rule flagged 363 of 14{,}400 numeric values (2.5\%) in 49 of the 60 columns; these were flagged, not removed. The most
affected column is \texttt{T4\_beta\_power} (46 trials), followed by \texttt{C4\_alpha\_power} (19 trials). Among the
categorical columns, \texttt{response\_accuracy} was missing for 2 trials (labelled ``Unknown'', consistent with the
ambiguous response markers), and \texttt{quality\_issues} was empty for 232 rows, which means no issue was recorded for
those trials (8 trials carry an issue). Full counts are in \texttt{reports/features\_cleaning\_log.csv}. The scaled table
(\texttt{features\_scaled\_demo.csv}) is fitted on the whole file for demonstration only; evaluation will refit on
training data only.

%===========================================================
\section{Phase 2: Feature Strategy}

\subsection{Spectral Features}
For each trial, the delay-to-probe window is extracted for the 19 EEG channels, and only that segment is read from disk
(loading entire recordings caused a memory error). Power spectral density is estimated with Welch's method
\cite{welch1967} using MNE-Python \cite{mne2013}, and band power is obtained by trapezoidal integration over theta
(4--8\,Hz), alpha (8--13\,Hz) and beta (13--30\,Hz). This gives $19\times3=57$ features per trial, stored in
\texttt{reports/eeg\_features.csv} (columns such as \texttt{Fp1\_theta\_power}) together with participant, trial,
condition, timing, response accuracy, trial quality and amplitude-screening columns. Theta and alpha power are common
markers of working-memory load, which motivates this first set. The features are \textbf{provisional}: they were extracted
before filtering, referencing and artifact rules were finalized.

\subsection{Scaling, Normalization and Encoding}
\begin{itemize}[leftmargin=*]
  \item \textbf{Scaling.} Band power is right-skewed and depends on each participant's overall amplitude. The log$_{10}$
  histograms in the band-power histograms in the EDA section are much more symmetric than raw power would be (theta still keeps a right tail), so the
  proposed approach for the first experiments is a log$_{10}$ transform followed by train-only standardization. Relative band power
  remains an option once filtering is settled. Whatever is adopted will be documented.
  \item \textbf{Normalization.} Standardization is fitted on training data only and applied to validation/test data to
  avoid leakage. In federated comparisons the scaling strategy (local per client or global) will be stated and kept
  identical across baselines.
  \item \textbf{Encoding.} Condition labels (5R--7M) are categorical and will be integer- or one-hot-encoded depending
  on the final target definition (six classes, or load level $\times$ task type). Response labels and quality flags are
  kept as separate columns.
\end{itemize}

%===========================================================
\section{Exploratory Data Analysis}
The EDA is descriptive; with two participants no inferential claims are made.

\safefig{trial_counts_by_condition.png}{0.8}{Trial counts by condition and participant (20 trials per condition per participant).}
\safefig{response_audit_status.png}{0.7}{Response/probe audit status by participant (structural review status, not behavioural accuracy).}
Of the 240 probes, 225 passed the structural audit and 15 need review (8 in sub-001, 7 in sub-002 in the figure).

\safefig{response_delay_distribution.png}{0.8}{Observed response-marker delays after probe onset.}
Most response markers arrive between about 1 and 3\,s after the probe, with a long right tail; a few delays exceed 5\,s,
which is what the five-second review heuristic targets.

\safefig{sub-001_representative_signals.png}{0.95}{Representative signal segments, sub-001.}
\safefig{sub-002_representative_signals.png}{0.95}{Representative signal segments, sub-002.}

\safefig{amplitude_by_participant.png}{0.6}{Maximum channel peak-to-peak amplitude per trial window by participant (log scale), with the 500\,$\mu$V screening heuristic.}
The two participants barely overlap: sub-001 windows centre near 410\,$\mu$V with a few high outliers (about 530, 700 and
1600\,$\mu$V), whereas every sub-002 window lies between about 2450 and 3000\,$\mu$V, far above the heuristic. This
separation is larger than any variation between conditions, so amplitude differences between recordings must be understood
before the features are pooled across participants.

\safefig{band_power_distributions.png}{0.95}{Distribution of theta, alpha and beta band power (log$_{10}$) over all trials and channels.}
Theta power is right-skewed with a long upper tail, alpha is close to symmetric, and beta shows a smaller secondary
bump around $\log_{10}\approx1.6$--$1.8$. The cleaning log flags \texttt{T4\_beta\_power} for 46 trials, so a single
channel may contribute to that bump, but this has not been checked.

\safefig{band_power_by_condition.png}{0.8}{Mean log$_{10}$ band power by condition (both participants pooled).}
Differences between conditions are small (means roughly 0.96--1.17 on the log scale) and show no clear ordering by set size
or task type. Because the participants differ strongly in amplitude, this pooled view is only descriptive, and no
condition effect should be read from it.

\safefig{band_power_by_participant.png}{0.8}{Mean log$_{10}$ band power by participant (descriptive, two participants).}
% TODO (team): add one sentence describing this figure after inspecting it.

\safefig{feature_correlation_heatmap.png}{0.85}{Correlation between the 57 log$_{10}$ band-power features.}
Features are ordered channel by channel (theta, alpha, beta), which produces the repeating three-column pattern. Features of
the same band are often strongly positively correlated across channels, and some frontal theta and beta features appear
negatively correlated, so the 57 features are far from independent. Dimensionality reduction or regularization should be
considered at the modeling stage.

%===========================================================
\section{Limitations and Remaining Work}
\begin{enumerate}[leftmargin=*]
  \item Two participants only: no generalization or cross-subject conclusions.
  \item Response-label mapping is provisional.
  \item Amplitude thresholds are heuristic, and the sub-001 boundary anomaly and the sub-002 flags are unexplained.
  \item Filtering, referencing and artifact handling are not finalized, so features may change; sub-002 windows are about six times larger in peak-to-peak range than sub-001 windows, and the cause is unknown.
  \item Next: finish the quality decisions, validate labels, expand participants, then run baseline, centralized,
  federated and personalized experiments with leakage-safe splits and suitable metrics (e.g.\ macro-F1).
\end{enumerate}

\section{Reproducibility}
Code is in \texttt{data\_scripts/} of the GitHub repository
(\url{https://github.com/EmanAdil122/federated-eeg-working-memory}); the README lists setup and execution order. Raw
EEG is not committed.

\begin{thebibliography}{9}
\bibitem{welch1967} P.~Welch, ``The use of fast Fourier transform for the estimation of power spectra,'' \emph{IEEE Trans.\ Audio Electroacoust.}, 15(2):70--73, 1967.
\bibitem{mne2013} A.~Gramfort et al., ``MEG and EEG data analysis with MNE-Python,'' \emph{Front.\ Neurosci.}, 7:267, 2013.
\bibitem{mcmahan2017} B.~McMahan et al., ``Communication-efficient learning of deep networks from decentralized data,'' \emph{AISTATS}, 2017.
\end{thebibliography}
\end{document}
